\documentclass[10pt,a4paper]{amsart}
\usepackage[margin=19mm]{geometry}
\usepackage{amsmath,amssymb,mathtools}
\usepackage[scr=rsfs,cal=euler]{mathalfa}
\usepackage{xfrac}
\usepackage[unicode,hidelinks,bookmarks=false]{hyperref}
\newcommand{\ie}{i.e.\ }
\newcommand{\cf}{cf.\ }
\newcommand{\resp}{respectively\ }
\newcommand{\Subsection}{\subsection}
\providecommand{\nobreakdash}{\nobreak\hbox{-}}
\setlength{\emergencystretch}{2em}
\multlinegap0pt
\usepackage{bm}
\usepackage{amssymb}
\usepackage{dsfont}
\newtheorem{lem}{Lemma}
\newcommand{\K}{\mathds{k}}
\newcommand{\Pp}{\mathcal{P}}
\newcommand{\Z}{\mathbb{Z}}
\title{Bicrossed products of generalized Taft algebras and Radford algebras}
\author{Rongchuan Xiong}
\date{Source: arXiv:2608.29536v1 (CC BY 4.0)}
\begin{document}
\maketitle

\noindent\emph{Source and licence.} Bicrossed products of generalized Taft algebras and Radford algebras. Version arXiv:2608.29536v1 (CC BY 4.0), distributed by arXiv under the Creative Commons Attribution 4.0 International licence, http://creativecommons.org/licenses/by/4.0/, as recorded in the arXiv licence metadata for this exact version and shown on the abstract page https://arxiv.org/abs/2608.29536v1. Full licence notice: \texttt{licenses/arxiv-2608.29536-CC-BY-4.0.txt}.\par\smallskip

\noindent\textit{Editorial scope.} Counted unit: the article's lem lem:primitive-R (source0.tex lines 460--475). The article's complete original proof of it is delivered with it (source0.tex lines 476--544, one \texttt{proof} environment of 69 lines). No mathematics is written, repaired or re-derived: the statement and the proof are the article's own lines, in the article's own notation, and no retained line is substituted or re-flowed.  No further source-local context is needed: every cross-reference of every retained line resolves inside the two delivered spans.  Nothing was omitted from any retained span, so the package carries no omission record.  No retained line contains a citation, so the wrapper carries no bibliography and no bibliography entry.  No retained line contains a figure environment, an image-inclusion command or drawing code, so no image is delivered and the package declares no asset dependency.  Editorial formatting only, in addition: an AMS \texttt{amsart} preamble that restates the article's own declarations and macros used by the retained lines, namely the environments \texttt{lem} on separate counters, together with the standard packages those lines need.  The retained environments print in this document as Lemma 1 in order of appearance, whereas the article prints the same items with the numbering of its own full document.  The complete native source of the article as distributed in the arXiv e-print for this exact version is bundled unchanged as \texttt{sources/208862/source0.tex}.
	\begin{lem}\label{lem:primitive-R}
		For every $s\in\Z/p\Z$,
		\[
		\Pp_{1,g^s}(R)=
		\begin{cases}
			0,&s=0,\\
			\K(g-1)\oplus\K x,&s=1,\\
			\K(g^s-1),&s\ne0,1.
		\end{cases}
		\]
		Moreover,
		\[
		\Pp_{g^s,g^s}(R)=0
		\qquad\text{for every }s\in\Z/p\Z.
		\]
	\end{lem}

	\begin{proof}
		Give $g$ degree $0$ and $x$ degree $1$. The relation
		\[
		xg=gx+(g^2-g)
		\]
		shows that commuting $x$ past a power of $g$ preserves the leading
		$x$-degree and creates only lower-degree correction terms. Consequently, for
		$0\le m<p$,
		\begin{equation}\label{eq:radford-leading}
			\Delta(g^ix^m)
			=\sum_{k=0}^{m}\binom{m}{k}
			g^{i+k}x^{m-k}\otimes g^ix^k
			+\{\text{terms of total $x$-degree $<m$}\}.
		\end{equation}
		
		The extra terms have strictly smaller total $x$-degree and cannot cancel
		the leading tensors in \eqref{eq:radford-leading}.
		
		Let
		\[
		z=\sum_{i,j=0}^{p-1}a_{i,j}g^ix^j\in\Pp_{1,g^s}(R).
		\]
		If $m\ge2$ is the maximal $x$-degree occurring in $z$, then the $k=1$ term
		in \eqref{eq:radford-leading} contains
		\[
		m\,a_{i,m}\,g^{i+1}x^{m-1}\otimes g^ix.
		\]
		Since $1\le m\le p-1$, the scalar $m$ is nonzero in $\K$. The right-hand
		side $z\otimes1+g^s\otimes z$ contains no term in which both tensor factors
		have positive $x$-degree. In the associated graded tensor coalgebra, the
		bidegree-$(m-1,1)$ tensor
		$g^{i+1}x^{m-1}\otimes g^ix$ can arise only from the $k=1$ term attached to
		$g^ix^m$; lower filtered-degree corrections cannot contribute to this
		component. Linear independence of the PBW tensor basis therefore gives
		$a_{i,m}=0$ for every $i$, a contradiction. Descending on $m$ eliminates
		all terms of $x$-degree at least two. If $p=2$, there are no
		PBW terms with degree $2\le j<p$, so the conclusion follows directly.
		
		Hence
		\[
		z=\sum_i a_i g^ix+\sum_i b_i g^i.
		\]
		Using
		\[
		\Delta(g^ix)=g^ix\otimes g^i+g^{i+1}\otimes g^ix,
		\]
		comparison of terms with $x$ in the first tensor factor gives $a_i=0$ for
		$i\ne0$. If $a_0\ne0$, comparison of the terms with $x$ in the second tensor
		factor gives $g=g^s$, hence $s=1$. Therefore the only possible $x$-term is a
		multiple of $x$, and it occurs only when $s=1$.
		
		For the remaining group-algebra part, comparison in the basis
		$\{g^i\otimes g^j\}$ yields
		\[
		\Pp_{1,1}(\K\langle g\rangle)=0,
		\qquad
		\Pp_{1,g^s}(\K\langle g\rangle)=\K(g^s-1)
		\quad(s\ne0).
		\]
		This proves the displayed formula for $\Pp_{1,g^s}(R)$. In particular,
		\[
		\Pp_{1,1}(R)=0.
		\]
		Finally, if $z\in\Pp_{g^s,g^s}(R)$, then
		\[
		\Delta(g^{-s}z)=g^{-s}z\otimes1+1\otimes g^{-s}z,
		\]
		so $g^{-s}z\in\Pp_{1,1}(R)=0$. Hence $z=0$.
	\end{proof}

\end{document}
